%% ma: L10-B13-0104
%% lop: 10
%% chuong: 5
%% bai: 13
%% dang: 10.13.13
%% loai: DS
%% muc: [NB, TH, TH, TH]
%% dap_an: "SSSĐ"
%% nguon: "chuyen-de-cac-so-dac-trung-cua-mau-so-lieu-khong-ghep-nhom-toan-10-chuong-trinh-moi.pdf (GV. Phan Nhật Linh) (P03), Bài 02 Dạng 2 Phần II Câu 8"
%% kien_thuc: "Lớp 10 Bài 13 (số trung bình, trung vị, tứ phân vị, mốt)"
%% trang_thai: chinh-thuc
%% ghi_chu: "Tự duyệt theo yêu cầu giáo viên 22/10/2026."
\begin{ex}
Cho bảng số liệu điểm kiểm tra môn Toán của $20$ học sinh:
\begin{center}
\begin{tabular}{|c|c|c|c|c|c|c|c|c|}\hline
Điểm & $4$ & $5$ & $6$ & $7$ & $8$ & $9$ & $10$ & Cộng\\\hline
Số học sinh & $1$ & $2$ & $3$ & $4$ & $5$ & $4$ & $1$ & $20$\\\hline
\end{tabular}
\end{center}
Xét tính đúng sai của các khẳng định sau:
\choiceTF{Mốt của mẫu số liệu là $5$}{Điểm trung bình của nhóm học sinh là $5$}{Trung vị của mẫu số liệu là $7$}{\True Tứ phân vị thứ nhất của mẫu số liệu là $6$}
\loigiai{a) Sai: $M_o=8$ (tần số $5$).

b) Sai: $\overline x=\dfrac{146}{20}=7{,}3$.

c) Sai: $M_e=\dfrac{x_{10}+x_{11}}2=\dfrac{7+8}2=7{,}5$.

d) Đúng: $Q_1=\dfrac{x_5+x_6}2=\dfrac{6+6}2=6$.}
\end{ex}
